\chapter{系统需求分析}
\label{chap:requirement}

\section{可行性分析}

\subsection{技术可行性}

项目选用 LayaAir 3.x 与 TypeScript，已实现 ECS 核心、战斗 Demo、元菜单、跑图生成、技能装配 UI 等模块。引擎支持 2D 预制体、帧循环与键盘输入；浏览器支持 Web Worker。SyDRA 等研究表明，理解引擎子系统结构有助于降低维护成本\cite{ullmann2025sydra}；Cantrell 等验证了商业游戏引擎用于复杂物理仿真的可行性\cite{cantrell2018unity}。配表经同步脚本发布至运行目录，技术路线可行。

\subsection{经济与社会可行性}

作为毕业设计原型，主要成本为开发与测试人力。游戏为奇幻战斗题材，须控制暴力表现并提示合理游戏时长；跑图中的休息节点提供节奏缓冲。

\section{功能性需求}

主要功能需求如表~\ref{tab:func-req} 所示。

\begin{table}[htbp]
  \centering
  \caption{主要功能性需求一览}
  \label{tab:func-req}
  \begin{tabular}{p{3.2cm}p{9.5cm}}
    \toprule
    \textbf{模块} & \textbf{需求描述} \\
    \midrule
    ECS 核心 & 创建/销毁实体；按类型存取组件；系统分组更新 \\
    玩法实体 & 玩家/怪物标签；移动、视图同步；hp/maxHp 与 modifier \\
    技能战斗 & 多 effect 技能；索敌；冷却；子弹表驱动伤害与穿透 \\
    怪物系统 & 波次刷怪；追逐；接触伤害；离屏回收 \\
    进度系统 & 经验、升级、随机奖励 \\
    元游戏 & 开始/选关；进入战斗；生存计时胜利 \\
    跑图 & 三大关 DAG；节点类型；种子可复现 \\
    技能装配 UI & Tab 暂停；三装备栏；拖拽装配 \\
    配置 & JSON 表注册加载；缺失时默认值与日志 \\
    \bottomrule
  \end{tabular}
\end{table}

\subsection{用例：局内战斗}

\textbf{参与者：}玩家。\textbf{前置条件：}配表已加载，\texttt{SimpleEcsDemo.init()} 完成。\textbf{主成功场景：}（1）输入移动；（2）自动施法；（3）子弹碰撞、穿透/分裂/连锁；（4）击杀获经验并升级；（5）Tab 打开技能面板暂停并装配。\textbf{扩展：}玩家死亡显示重启面板。

\subsection{用例：元游戏进入战斗}

用户在 \texttt{MetaFlowController} 选择关卡或跑图节点，触发 \texttt{onEnterCombat}，显示战斗场景、初始化 Demo、启动生存计时（Boss 关使用更长胜利时间）。

\section{非功能性需求}

\textbf{性能：}目标 60 FPS；同屏怪物受常量与波次控制；配合对象池与 Worker。\textbf{可维护性：}模块边界清晰；静态配置集中于 \texttt{defines}。\textbf{兼容性：}优先 Chromium；Worker 不可用时主线程回退\cite{cantrell2018unity}。

\section{需求优先级}

\begin{table}[htbp]
  \centering
  \caption{需求优先级（MoSCoW）}
  \begin{tabular}{cl}
    \toprule
    \textbf{级别} & \textbf{需求} \\
    \midrule
    Must & ECS 战斗、子弹碰撞、移动、刷怪、配表 \\
    Must & 主菜单进战斗、生存胜利计时 \\
    Should & 跑图、技能装配、升级奖励、对象池与 Worker \\
    Could & 法杖三槽完整 UI、全量 Effect 独立特效 \\
    Won't & 联机、账号、支付（本期） \\
    \bottomrule
  \end{tabular}
\end{table}

\section{运行环境}

建议配置：四核 CPU、8GB 内存、支持 WebGL2 的浏览器；开发环境为 LayaAir IDE 3.x、TypeScript、Node.js 工具链。

\section{本章小结}

本章论证了可行性，列出功能与非功能需求及典型用例，下一章给出总体设计。
